\documentclass[5pt, a4paper]{article}

\usepackage[utf8]{inputenc}
\usepackage[T1]{fontenc}
\usepackage{textcomp}
\usepackage{amsmath, amssymb}
\usepackage{multicol}
\usepackage{proof}
\usepackage{enumitem}
\usepackage{fancyhdr}
\usepackage[margin=1in]{geometry}
% figure support
\usepackage{import}
\usepackage{xifthen}
\pdfminorversion=7

\usepackage[most]{tcolorbox}
\usepackage{pdfpages}
\usepackage{transparent}
\newcommand{\incfig}[1]{%
    \def\svgwidth{\columnwidth}
    \import{./figures/}{#1.pdf_tex}
}

\tcbset{
colback=gray!5,
colframe=blue!50!white,
coltitle=black,
arc=0pt,
boxrule=-0.1mm,
fonttitle=\bfseries,
coltitle=black,
left=2.5mm,
leftrule=1mm,
right=3.5mm,
colbacktitle=black!10!white,
toptitle=0.75mm,
bottomtitle=0.5mm,
}
\begin{document}
\section{System of Linear Equation}
This chapter primarily consists of the application of matrices in solving system
of linear equations. Since a lot of this chapter is just review of linear
algebra, I 'll omit a lot of it, and focus on the application instead.
\subsection*{Matrix Operations}
In practice, matrix is simply a 2-dimension array, and operations on and between
matrices can be summed up as iterating over the array. Let's talk addition
first, which is performed on matrices of the same size. To compute it, we simply
need to iterate on each entries of the matrix and add a corresponding entry on
the other matrix. For a matrix $A$ of size $m \times n$, we would need to
perform $m*n$ operations.

Suppose matrices  $A_{m \times p}$ and $B_{p \times n}$. A multiplication
between them would involve, for each entry in the  new matrix of size $m \times
n$, an initialization followed by $2p - 2$ operations (addition and
multiplication). Thus, we have a total of $nm(2p-1)$ total operations.

However, we can further optimize this for matrices with special structures.
Consider an upper triangular matrix $B_{n \times n}$ and some matrix $A_{m
\times n}$. We can further minimize the operations into $mn\frac{n+1}{2}$
additions, and $mn\frac{\left(  n-1 \right) }{2}$ multiplications.

Furthermore, for a multiplication between two lower triangular matrix, the total
operations would be $n(n+1)(n+2)/6$ multiplication and $(n-1)n(n+1) / 6$
additions.

\subsection*{Norms and Errors}
There are a couple of norms that each can be used to measure the `size` of a
vector, and consequently the error of matrix operations. This course would cover
the euclidiean norm (2-norm), the p-norm and the infinity norm.

Let's begin with the general case of p-norm. It is defiend by  \[
\|x\|_p = \left(\sum_{i = 1}^{n} |x_{i}|^p  \right)^{\frac{1}{p}}
.\] 
This way, you can form the euclidian norm \[
\|x\|_2 = \left(\sum_{i = 1}^{n} |x_{i}|^2  \right)^{\frac{1}{2}}
.\] 
and the infinity norm, which is the limit as $p \to \infty$. It yields \[
    \|x\|_{\infty} = \max_i |x_i|
,\] 
as in, it is defined by the largest entry. Generalizing it to matrices, we can
defien the following:
\[
    \|A\|_1 = \text{max absolute column sum} 
\]
\[
    \|A\|_{\infty} = \text{max absolute row sum}
\]
\[
    \|A\|_{2} = \text{max singular value in $A$}
.\] 

Now that we have a way of quantifying the size of a matrix, we can calculate the
condition number of a matrix. Recall the definition of the condition number: \[
\kappa = \frac{|\text{forward error}|}{|\text{backward error}|} = 
\frac{\left|\frac{y}{\Delta y}\right|}{\left|\frac{x}{\Delta x}\right|}
.\] 
Rewriting it in terms of the transformation $A(x) = b$, \[
\kappa =  \frac{\|\Delta x\| \|b\|}{\|\Delta b\| \|x\|}
= \frac{\|A^{-1}(\Delta b)\| \|A(x)\|}{\|\Delta b\| \|x\|}
=\frac{\|A^{-1}\|\|\Delta b\| \|A\|\|x\|}{\|\Delta b\| \|x\|}
.\] 
\[
\kappa \le \|A^{-1}\|\|A\|
.\] 

\subsection*{LU Factorization}
For an LU factorization, we have two main algorithms: 
\begin{enumerate}
    \item Gaussian  Elimination.
    \item Cholseky Algorithm
\end{enumerate}
In general, Cholesky is the one that is preferred, as it is a bit more
effecient and stable. Regardless, Gaussian Elimination is worth knowing.

Gaussian Elimination is the method most people do by hand. It involves forming
an upper triangular matrix, while storing the row operations in a matrix form.
The resulting upper triangualr matrix is $U$, whilst the inverse of the row
operation matrix is $L$. Writing the algorithm for a Gaussian Elimination is
also quite straightforward: simple use forward subtitution to form the upper
triangular matrix, whilst accumulating the operations in a seperate $L$ vector
of the same size.

One might notice that when the pivot entry of a matrixis zero, or very  close to
zero, the algorithm might run into some problem. To remedy this, we turn to
pivoting, as in swapping the rows of a matrix such that zero divisions is
avoided. The standard way to do this is that, for every column of a matrix, we
pick the row that contains the largest entry below the diagonal, and swap that
row as the pivot row for that iteration. This adds extra complexity in the
accumulation of $L$, since it now has to include pivoting matrices.

\subsection*{Cholesky Factorization}
Cholesky Factorization can only be done on matrices that is positive definite
and symmetric. It is defined as a matrix that has a nonnegative determinant on
all it'sleading principal submatrices.

The idea of a Cholseky Factorization can be demostratedd as such: Suppose a
symetric positive definite matrix $A_{3\times 3}$. It's LU facotrization is can be written
as \[
    A = \begin{pmatrix} 
    l_{11} & 0 & 0 \\ 
    l_{21} & l_{22} & 0 \\
    l_{31} & l_{32} & l_{33} \\
\end{pmatrix} 
\begin{pmatrix} 
    l_{11} & l_{21} & l_{31} \\ 
    0 & l_{22} & l_{32} \\
    0 & 0 &l_{33}    
\end{pmatrix} 
.\] 
We can compute each $l_{ij}$ based on the entry of $A$. Specifically, we can
write the following: 
\begin{align*}
    &a_{11} = l_{11}^2 \\
    &a_{21} = l_{11}l_{21} \\
    & a_{31} = l_{11}l_{31} \\
    & a_{22} = l_{21}l_{22} \\
    & a_{32} = l_{21}l_{31} + l_{22}l_{32} \\
    &a_{33} = l_{31}^2 + l_{32}^2 + l_{33}^2 
.\end{align*}
We'd solve for each value of $l_{ij}$, which can be done following the above
sequence.

\end{document}
